\section{Translation from controlled prose}
\label{sec:prop-syntax-semantics-03}

A translation requires a key assigning one propositional variable to each
atomic statement.  The key is part of the translation; a formula without its
key does not preserve the content of the prose.

\begin{example}[Necessary and sufficient conditions]
Use the key
\[
  P:\text{the archive is readable},
  \qquad Q:\text{the checksum is valid}.
\]
Then:
\begin{center}
\begin{tabular}{@{}ll@{}}
\toprule
Controlled prose & Formula \\ \midrule
If the archive is readable, then the checksum is valid. & \(P\to Q\) \\
The archive is readable only if the checksum is valid. & \(P\to Q\) \\
The archive is readable if the checksum is valid. & \(Q\to P\) \\
The archive is readable iff the checksum is valid. & \(P\leftrightarrow Q\) \\
The archive is readable unless the checksum is invalid. & \(\neg\neg Q\to P\) \\
\bottomrule
\end{tabular}
\end{center}
The word ``only'' marks a necessary condition, and it reverses the effect
of a bare ``if'': the ``if'' of the third row introduces the antecedent,
while the ``only if'' of the second row introduces the consequent, so the
main clause of the second row supplies the antecedent and the first two
rows share one formula.
For the final row, ``\(P\) unless \(R\)'' is \(\neg R\to P\).  Here
\(R\) is ``the checksum is invalid,'' namely \(\neg Q\), so the translation
is \(\neg\neg Q\to P\), and that formula is what the table records.
Simplifying it to \(Q\to P\) is a further step, justified by an equivalence
of classical propositional logic and not part of the translation.
\end{example}

\begin{example}[Exclusive disjunction]
``Exactly one of \(P\) and \(Q\)'' is translated by
\[
  (P\lor Q)\land\neg(P\land Q).
\]
The formula \(P\lor Q\) alone is inclusive: it is true when both disjuncts are
true.
\end{example}

